\section{Discussion}
\label{sec:discussion}

\subsection{Interpretation}

The central observation of this study concerns how two anatomical priors behave
when introduced into the same fine-grained grading backbone. Supervising
anatomical position alone, or weak lesion extent alone, left the composite
grading risk worse than the unmodified baseline. Supplying both together did not
inherit that degradation, and produced the best observed composite risk of the
four configurations. This is a non-additive empirical pattern: the behavior of
the joint configuration is not what either single-prior configuration would
suggest on its own.

We can offer a plausible account, while noting that our experiments do not
isolate the mechanism. Each prior, taken alone, adds a training objective whose
optimum need not align with fine-grained grading. Position classification
rewards features that discriminate scan location, and weak-box attention rewards
features that concentrate on annotated regions; neither objective is
intrinsically about septal texture. Injected in isolation, such a signal may
displace grading-relevant capacity. When both are present, the position estimate
supplies the anatomical frame within which the lesion attention is interpreted,
and the two together may describe evidence more completely than either does
alone. This account is consistent with the observations but is not established
by them, and disentangling multi-task regularization from prediction injection
would require a configuration that retains both auxiliary losses while
suppressing the residual injection. We did not run that configuration.

A second observation is more directly actionable. The joint configuration
produced two additional outputs -- a six-class anatomical position estimate and
a weak lesion localization map -- for a $5.9\%$ parameter and $2.8\%$ arithmetic
overhead, with median single-image latency remaining below $4$~ms. Both outputs
are derived from the image alone and require no metadata at inference. In a
screening workflow, a position estimate permits automatic checking that the
standardized six-position protocol was followed, and a localization map gives
the reader something to inspect alongside a numeric score. Neither capability
existed in the baseline, and neither required additional annotation beyond what
the cohort already carried.

The stage-wise analysis argues against a single summary of where difficulty
lies. Mean absolute error was largest at F2 for both configurations, whereas
recall was lowest at F1. These two metrics therefore identify different stages
as hardest. It is clinically plausible that mild fibrosis presents subtler and
more heterogeneous sonographic
manifestations than either normal parenchyma or advanced disease, which would be
consistent with the observed F1 recall; we note this as a plausible reading
rather than a demonstrated explanation.

Finally, the comparison between the two resize conventions favoured anisotropic
stretching over aspect-preserving letterbox padding on the composite endpoint.
We retained stretching as the default preprocessing for all modeling
configurations on this basis, and report the letterbox control in
Table~\ref{tab:main} so that the choice is visible rather than implicit.

\subsection{Limitations}

Several limitations bound what the present study establishes.

\textbf{Etiological and institutional scope.} The cohort consists exclusively of
\emph{Schistosoma japonicum}-associated liver fibrosis. The septal, network-like
phenotype that makes anatomical context plausible here differs from the diffuse
parenchymal change of viral, alcoholic and metabolic liver disease, and none of
our results speak to those etiologies. Separately, although the cohort is
multicenter and the test split is patient-disjoint, the same centers were
represented during training and testing. The present study therefore does not
establish generalization to unseen institutions or acquisition environments,
which would require dedicated external validation at centers absent from
training.

\textbf{Test cohort composition and center sizes.} The test cohort contained
60~patients in each of the four stages, whereas training and validation followed
the cohort distribution. Aggregate test performance should accordingly not be
read as a prevalence-weighted estimate for a screening population. The four test
centers also contributed very unequal numbers of patients, from 2 to 115, so
center-specific estimates -- reported in the Supplementary Material -- are
descriptive, and those from the smallest centers should be interpreted with
caution.

\textbf{Single configuration and single seed.} The comparison rests on one
prespecified seed, and we do not present it as evidence about run-to-run
variability; the margin separating the joint configuration from the baseline on
the primary endpoint is small relative to what a single run can resolve, and we
therefore describe the joint configuration as the best observed rather than as an
established improvement. The design was also evaluated with the prespecified
ResNet-50-based framework only, and its robustness across alternative backbone
architectures remains to be established.

\textbf{Reference standard and lesion supervision.} Grading labels derive from
expert ultrasonographic reading rather than histopathology, as is necessary for a
non-invasive screening cohort of this scale, so the model is calibrated to expert
consensus rather than to tissue-level ground truth. The lesion branch is
supervised by coarse bounding boxes; the reported Dice and IoU quantify agreement
with those weak annotations and are not measurements of lesion segmentation
accuracy, for which no pixel-level reference exists in this cohort.

\section{Conclusion}
\label{sec:conclusion}

We presented SynAP-Fib, a fine-grained ultrasound grading model for
schistosomiasis-associated liver fibrosis in which anatomical position and weak
lesion evidence serve as training-time supervision and as inference-time context
through gated, gradient-isolated residual injection, requiring no anatomical
metadata at deployment. A $2\times2$ comparison of the two priors revealed a
non-additive empirical pattern: either prior alone degraded the prespecified
composite grading risk, whereas their combination did not and attained the best
observed composite risk among the four configurations. The joint configuration
additionally emitted anatomical position and weak lesion localization outputs at
a small computational cost while retaining baseline-level grading performance.
Establishing that these gains hold across seeds, backbones and institutions
absent from training remains future work, as does isolating whether the benefit
of the joint configuration arises from auxiliary-task regularization or from the
injection of the predicted anatomical context itself.
